\chapter{Guide de reproduction de l'environnement}
\label{ann:reproduction}

Ce guide permet à un tiers de reproduire intégralement l'environnement de développement
et d'obtenir les résultats rapportés dans ce mémoire.

\section*{H.1 — Versions exactes des dépendances}

\begin{table}[H]
\centering
\caption{Versions exactes des dépendances principales du projet}
\label{tab:versions}
\begin{tabular}{llll}
\toprule
\textbf{Composant} & \textbf{Version (min)} & \textbf{Composant} & \textbf{Version (min)} \\
\midrule
Python         & 3.11.x            & FastAPI            & $\geq$0.100.0 \\
Node.js        & 20 LTS            & SQLAlchemy         & $\geq$2.0.36 \\
Docker         & 25.x              & Pydantic           & $\geq$2.10.0 \\
Docker Compose & 2.24.x            & Ultralytics (YOLO) & $\geq$8.1.0 \\
PostgreSQL     & 16-alpine         & React              & 18.2.0 \\
Ollama         & $\geq$0.1.7       & Vite               & 5.1.5 \\
MLflow         & $\geq$2.10.0      & ChromaDB           & $\geq$0.4.24 \\
DVC            & $\geq$3.45.0      & sentence-transformers & $\geq$2.5.0 \\
GeoAlchemy2    & $\geq$0.14.7      & torch              & $\geq$2.2.0 \\
paho-mqtt      & $\geq$1.6.1       & Three.js           & 0.183.2 \\
\bottomrule
\end{tabular}
\end{table}

\section*{H.2 — Procédure de démarrage}

\begin{lstlisting}[language=bash,
  caption={Script de démarrage complet — Linux/macOS/WSL2},
  label={lst:setup}]
# 1. Cloner le repository
git clone https://github.com/<user>/smart-farm-ai.git
cd smart-farm-ai

# 2. Configurer les variables d'environnement
cp backend/.env.example backend/.env
# Editer backend/.env : DATABASE_URL, GROQ_API_KEY, META_WHATSAPP_TOKEN

# 3. Demarrer les services Docker (db + backend + frontend actifs)
# docker-compose.yml :
#   db       : postgres:16-alpine  (port 5432, volume postgres_data)
#   backend  : Dockerfile Python   (port 8000, uvicorn --workers 2)
#   frontend : Dockerfile nginx    (port 80, proxy /api/v1 -> backend)
#   caddy    : caddy:2-alpine      (commenté - activer pour HTTPS)
docker compose --env-file .env.production up -d --build

# 4. Initialiser la base de donnees
docker exec smart_farm_backend python -m app.scripts.seed_db

# 5. Telecharger le modele LLM local (Labess-7B, ~4.7 Go)
docker exec smart_farm_backend ollama pull labess:7b

# 6. Lancer la suite de tests avec couverture
docker exec smart_farm_backend pytest --cov=app --cov-report=html -v

# 7. Acces aux services
# Frontend  : http://localhost:3000
# API docs  : http://localhost:8000/docs   (Swagger UI)
# MLflow UI : http://localhost:5000
# Adminer   : http://localhost:8080
\end{lstlisting}

\section*{H.3 — Reproduction des expériences ML}

\begin{lstlisting}[language=bash,
  caption={Reproduction des entraînements YOLO et ML avicoles (3 runs + 5-fold CV)},
  label={lst:ml_repro}]
# Reproduction des 3 runs YOLO independants (seeds documentes)
for seed in 42 137 255; do
  python backend/app/scripts/train_yolo.py \
    --data data/yolo_dataset.yaml \
    --model yolov8n.pt \
    --epochs 100 --batch 16 --imgsz 640 \
    --seed $seed \
    --project mlruns/yolo_runs \
    --name "run_seed_${seed}"
done
# Resultat attendu : mAP@50 moyen = 84.7% +/- 0.45%

# Validation croisee 5-fold modeles avicoles (TimeSeriesSplit)
python backend/app/scripts/train_poultry_ml.py \
  --data data/poultry_feed_logs.csv \
  --cv-folds 5 \
  --time-series-split \
  --mlflow-tracking-uri ./mlruns
# Resultat attendu : R2(FCR) = 0.891 +/- 0.024
\end{lstlisting}

